% ============================================================
% CAPITOLO 4 — Progetto 1: Hello World in 0-Code
% ============================================================
\chapter{Progetto 1: Hello World in 0-Code}\index{Python}

\section*{Cosa Costruirai}
Un programma Python interattivo che:
\begin{itemize}
  \item Saluta l'utente per nome
  \item Calcola l'età a partire dall'anno di nascita
  \item Genera una ``fortuna del giorno'' casuale
  \item Ha test automatici funzionanti
\end{itemize}

\textbf{Tempo stimato}: 15--20 minuti

% ---------------------------------------------------------
\section{Preparare il Progetto (ADLC Fase 0-2)}

\begin{praticabox}[title={Pratica --- Setup iniziale}]
\begin{enumerate}
  \item Apri VS~Code
  \item Crea una nuova cartella: \texttt{hello-0code}
  \item Apri la cartella in VS~Code (\texttt{File $\rightarrow$ Open Folder})
  \item Crea un file \texttt{\_CONTEXT.md} nella root con questo contenuto:
\end{enumerate}
\end{praticabox}

\begin{lstlisting}[language={},title={\texttt{\_CONTEXT.md} per Hello 0-Code}]
# Progetto: Hello 0-Code

## Scopo
Un programma Python interattivo che saluta l'utente, calcola
la sua eta' e genera una "fortuna del giorno" casuale.

## Tecnologie
- Python 3.11+
- Solo librerie standard (random, datetime)
- Test con pytest

## Struttura
hello-0code/
+-- _CONTEXT.md
+-- main.py          <- Entry point
+-- greeter.py       <- Logica di saluto e calcolo
+-- fortunes.py      <- Lista di fortune e generatore
+-- tests/
    +-- test_greeter.py
    +-- test_fortunes.py

## Convenzioni
- snake_case per tutto
- Type hints su tutte le funzioni
- Docstring per ogni funzione pubblica
- Print formattato con f-strings

## Vincoli
- Solo librerie standard Python. Nessun pip install.
- L'anno di nascita deve essere validato (1900-anno corrente)
- Il nome non puo' essere vuoto
- Gestisci input non validi con messaggi chiari
\end{lstlisting}

\begin{checkpointbox}
Dovresti avere una cartella con un solo file: \texttt{\_CONTEXT.md}. Tutto il resto lo genererà l'IA.
\end{checkpointbox}

% ---------------------------------------------------------
\section{La Prima Richiesta (ADLC Fase 3-4)}

\begin{praticabox}[title={Pratica --- Generare il codice}]
\begin{enumerate}
  \item Apri Copilot Chat in Agent Mode (\texttt{Ctrl+Shift+I}, seleziona ``Agent'')
  \item Scrivi questa richiesta:
\end{enumerate}
\begin{lstlisting}[language={}]
Leggi il file _CONTEXT.md e implementa l'intero progetto
seguendo la struttura e i vincoli definiti.

Il programma deve:
1. Chiedere il nome dell'utente
2. Chiedere l'anno di nascita
3. Calcolare l'eta'
4. Mostrare un saluto personalizzato con l'eta'
5. Mostrare una fortuna del giorno casuale (10+ fortune)

Crea anche i test unitari per greeter.py e fortunes.py.
\end{lstlisting}
\begin{enumerate}
  \setcounter{enumi}{2}
  \item \textbf{Osserva l'IA al lavoro}: Copilot in Agent Mode creerà i file uno per uno e ti chiederà se accettare le modifiche.
  \item \textbf{Accetta o correggi}: Se il codice è corretto, accetta. Se noti qualcosa di strano, chiedi una modifica.
\end{enumerate}
\end{praticabox}

% ---------------------------------------------------------
\section{Analisi dell'Output}

Vediamo cosa l'IA dovrebbe aver generato. Non preoccuparti se il tuo output è leggermente diverso --- l'importante è che rispetti la struttura del \texttt{\_CONTEXT.md}.

\subsection*{Output atteso: \texttt{greeter.py}}

\begin{lstlisting}[language=Python,title={\texttt{greeter.py}}]
from datetime import datetime


def get_user_name() -> str:
    """Chiede il nome all'utente e lo valida."""
    while True:
        name = input("Come ti chiami? ").strip()
        if name:
            return name
        print("Il nome non puo' essere vuoto. Riprova.")


def get_birth_year() -> int:
    """Chiede l'anno di nascita e lo valida."""
    current_year = datetime.now().year
    while True:
        try:
            year = int(input("In che anno sei nato/a? "))
            if 1900 <= year <= current_year:
                return year
            print(f"L'anno deve essere tra 1900 e {current_year}.")
        except ValueError:
            print("Inserisci un numero valido.")


def calculate_age(birth_year: int) -> int:
    """Calcola l'eta' a partire dall'anno di nascita."""
    return datetime.now().year - birth_year


def greet(name: str, age: int) -> str:
    """Genera il messaggio di saluto personalizzato."""
    return f"Ciao {name}! Hai {age} anni."
\end{lstlisting}

\subsection*{Cosa osservare}

Controlla che l'IA abbia rispettato il \texttt{\_CONTEXT.md}:

\begin{center}
\begin{tabularx}{\textwidth}{X c}
\toprule
\textbf{Vincolo del contesto} & \textbf{Rispettato?} \\
\midrule
snake\_case per funzioni & \checkmark \\
Type hints & \checkmark \\
Docstring & \checkmark \\
Validazione anno (1900--corrente) & \checkmark \\
Nome non vuoto & \checkmark \\
Solo librerie standard & \checkmark \\
\bottomrule
\end{tabularx}
\end{center}

\begin{tipbox}
Questa verifica è il tuo lavoro principale come sviluppatore 0-code. Non devi scrivere il codice, ma devi saper \textbf{leggere e validare} che rispetti i requisiti.
\end{tipbox}

% ---------------------------------------------------------
\section{Eseguire e Testare (ADLC Fase 5)}

\begin{praticabox}[title={Pratica --- Esecuzione del programma}]
Apri il terminale integrato e esegui:
\begin{lstlisting}[language=bash]
python main.py
\end{lstlisting}
Interagisci con il programma:
\begin{lstlisting}[language={}]
Come ti chiami? Marco
In che anno sei nato/a? 1990
Ciao Marco! Hai 36 anni.
La tua fortuna del giorno: "Il codice migliore e'
quello che non devi scrivere."
\end{lstlisting}
\end{praticabox}

\begin{praticabox}[title={Pratica --- Esecuzione dei test}]
\begin{lstlisting}[language=bash]
pip install pytest
python -m pytest tests/ -v
\end{lstlisting}
Output atteso:
\begin{lstlisting}[language={}]
tests/test_greeter.py::test_calculate_age PASSED
tests/test_greeter.py::test_greet PASSED
tests/test_greeter.py::test_greet_message_format PASSED
tests/test_fortunes.py::test_get_fortune_returns_string PASSED
tests/test_fortunes.py::test_fortune_not_empty PASSED
========================= 5 passed =========================
\end{lstlisting}
\end{praticabox}

\begin{checkpointbox}
Se il programma funziona e i test passano, hai completato il tuo primo progetto in 0-code.
\end{checkpointbox}

% ---------------------------------------------------------
\section{Iterare e Migliorare}

Il vero potere del 0-code emerge nell'iterazione.

\begin{praticabox}[title={Pratica --- Aggiungere una funzionalità}]
Nella chat di Copilot, scrivi:
\begin{lstlisting}[language={}]
Aggiungi una funzionalita' "zodiac": dopo il saluto, mostra
il segno zodiacale dell'utente basato sull'anno di nascita
(zodiaco cinese). Aggiungi la logica in un nuovo file
zodiac.py con i relativi test. Aggiorna main.py.
\end{lstlisting}
L'IA creerà \texttt{zodiac.py}, i test corrispondenti, e aggiornerà \texttt{main.py}.
\end{praticabox}

\begin{praticabox}[title={Pratica --- Correggere un bug}]
Se il programma ha un comportamento inatteso, non modificare il codice a mano. Descrivi il problema a Copilot:
\begin{lstlisting}[language={}]
Quando inserisco l'anno 2024 come anno di nascita, il
programma dice che ho 2 anni. Dovrebbe calcolare l'eta'
usando anche mese e giorno, non solo l'anno. Correggi.
\end{lstlisting}
\end{praticabox}

% ---------------------------------------------------------
\section{Lezioni dal Primo Progetto (ADLC Fase 6)}

\begin{enumerate}
  \item \textbf{Il \texttt{\_CONTEXT.md} guida tutto}: struttura dei file, convenzioni, vincoli --- l'IA li ha seguiti perché li hai definiti nel contesto.
  \item \textbf{La richiesta iniziale può essere generica}: i dettagli stanno nel contesto, la richiesta indica l'obiettivo.
  \item \textbf{L'iterazione è gratuita}: aggiungere funzionalità o correggere bug è veloce come descrivere il problema.
  \item \textbf{Tu revisioni, l'IA implementa}: il tuo lavoro è leggere il codice e verificare che rispetti i requisiti.
\end{enumerate}

\begin{praticabox}[title={Pratica --- Aggiorna il contesto}]
Se l'IA ha commesso errori ricorrenti, aggiungili come vincoli:
\begin{lstlisting}[language={}]
## Lezioni Apprese
- Calcolare l'eta' usando datetime.date completo
- I test devono usare fixture, non valori hardcoded
\end{lstlisting}
\end{praticabox}

% ---------------------------------------------------------
\section*{Riepilogo del Progetto}

\begin{center}
\begin{tabularx}{\textwidth}{l X c}
\toprule
\textbf{Fase ADLC} & \textbf{Cosa hai fatto} & \textbf{Tempo} \\
\midrule
Fase 0--2 & Creato \texttt{\_CONTEXT.md} & 5~min \\
Fase 3--4 & Generato codice con Copilot & 5~min \\
Fase 5 & Testato e verificato & 3~min \\
Iterazione & Aggiunto zodiaco e fix bug & 5~min \\
Fase 6 & Aggiornato contesto & 2~min \\
\midrule
\textbf{Totale} & \textbf{Applicazione completa con test} & \textbf{$\sim$20~min} \\
\bottomrule
\end{tabularx}
\end{center}
